\subsection*{Addition on $\mathbb{Q}$-Classes}

\begin{definitionbox}{Definition}
\begin{definition}[Addition on $\mathbb{Q}$-Classes]
\label{def:addition-on-q-classes}
Define addition on representatives by
\[
[a,b]+[c,d]:=[a\cdot d+b\cdot c,\ b\cdot d].
\]
\end{definition}
\end{definitionbox}

\begin{dependencies}
\begin{itemize}
  \item \hyperref[def:rational-class-notation]{Rational Class Notation}
  \item \hyperref[def:addition-on-z-classes]{Addition on $\mathbb{Z}$}
  \item \hyperref[def:multiplication-on-z-classes]{Multiplication on $\mathbb{Z}$}
\end{itemize}
\end{dependencies}

\begin{lemmabox}{Lemma}
\begin{lemma}[Addition Lands in $\operatorname{Pre}\mathbb{Q}$]
\label{lem:addition-denominator-nonzero-q}
If $b\ne 0_{\mathbb{Z}}$ and $d\ne 0_{\mathbb{Z}}$, then
\[
b\cdot d\ne 0_{\mathbb{Z}}.
\]
\end{lemma}
\end{lemmabox}

\begin{dependencies}
\begin{itemize}
  \item \hyperref[thm:z-has-no-zero-divisors]{$\mathbb{Z}$ Has No Zero Divisors}
\end{itemize}
\end{dependencies}

\begin{lemmabox}{Lemma}
\begin{lemma}[Addition Respects Fraction Equivalence on the Left]
\label{lem:addition-on-q-respects-left-equivalence}
If $[a,b]=[a',b']$, then
\[
[a,b]+[c,d]=[a',b']+[c,d].
\]
\end{lemma}
\end{lemmabox}

\begin{dependencies}
\begin{itemize}
  \item \hyperref[def:addition-on-q-classes]{Addition on $\mathbb{Q}$}
  \item \hyperref[def:fraction-equivalence]{Fraction Equivalence}
  \item \hyperref[prop:equality-criterion-for-rational-classes]{Equality Criterion for Rational Classes}
  \item \hyperref[cor:multiplicative-cancellation-on-z]{Multiplicative Cancellation on $\mathbb{Z}$}
\end{itemize}
\end{dependencies}

\begin{lemmabox}{Lemma}
\begin{lemma}[Addition Respects Fraction Equivalence on the Right]
\label{lem:addition-on-q-respects-right-equivalence}
If $[c,d]=[c',d']$, then
\[
[a,b]+[c,d]=[a,b]+[c',d'].
\]
\end{lemma}
\end{lemmabox}

\begin{dependencies}
\begin{itemize}
  \item \hyperref[def:addition-on-q-classes]{Addition on $\mathbb{Q}$}
  \item \hyperref[def:fraction-equivalence]{Fraction Equivalence}
  \item \hyperref[prop:equality-criterion-for-rational-classes]{Equality Criterion for Rational Classes}
  \item \hyperref[cor:multiplicative-cancellation-on-z]{Multiplicative Cancellation on $\mathbb{Z}$}
\end{itemize}
\end{dependencies}

\begin{lemmabox}{Lemma}
\begin{lemma}[Addition Respects Fraction Equivalence]
\label{lem:addition-on-q-respects-equivalence}
If $[a,b]=[a',b']$ and $[c,d]=[c',d']$, then
\[
[a,b]+[c,d]=[a',b']+[c',d'].
\]
\end{lemma}
\end{lemmabox}

\begin{dependencies}
\begin{itemize}
  \item \hyperref[def:addition-on-q-classes]{Addition on $\mathbb{Q}$}
  \item \hyperref[lem:addition-on-q-respects-left-equivalence]{Addition Respects Fraction Equivalence on the Left}
  \item \hyperref[lem:addition-on-q-respects-right-equivalence]{Addition Respects Fraction Equivalence on the Right}
\end{itemize}
\end{dependencies}

\begin{theorembox}{Theorem}
\begin{theorem}[Addition on $\mathbb{Q}$ Is Well-Defined]
\label{thm:addition-on-q-well-defined}
Addition determines a well-defined binary operation
$\mathbb{Q}\times\mathbb{Q}\to\mathbb{Q}$.
\end{theorem}
\end{theorembox}

\begin{dependencies}
\begin{itemize}
  \item \hyperref[def:addition-on-q-classes]{Addition on $\mathbb{Q}$}
  \item \hyperref[lem:addition-on-q-respects-equivalence]{Addition Respects Fraction Equivalence}
\end{itemize}
\end{dependencies}

\begin{theorembox}{Theorem}
\begin{theorem}[Addition on $\mathbb{Q}$ Is Associative]
\label{thm:addition-on-q-associative}
For all $x,y,z\in\mathbb{Q}$,
\[
(x+y)+z=x+(y+z).
\]
\end{theorem}
\end{theorembox}

\begin{theorembox}{Theorem}
\begin{theorem}[Addition on $\mathbb{Q}$ Is Commutative]
\label{thm:addition-on-q-commutative}
For all $x,y\in\mathbb{Q}$,
\[
x+y=y+x.
\]
\end{theorem}
\end{theorembox}

\begin{dependencies}
\begin{itemize}
  \item \hyperref[def:addition-on-q-classes]{Addition on $\mathbb{Q}$}
  \item \hyperref[thm:addition-on-q-well-defined]{Addition on $\mathbb{Q}$ Is Well-Defined}
  \item \hyperref[thm:addition-on-z-commutative]{Addition on $\mathbb{Z}$ Is Commutative}
\end{itemize}
\end{dependencies}

\begin{theorembox}{Theorem}
\begin{theorem}[Zero Is a Left Additive Identity on $\mathbb{Q}$]
\label{thm:zero-left-additive-identity-on-q}
For every $x\in\mathbb{Q}$,
\[
0_{\mathbb{Q}}+x=x.
\]
\end{theorem}
\end{theorembox}

\begin{corollarybox}{Corollary}
\begin{corollary}[Zero Is a Right Additive Identity on $\mathbb{Q}$]
\label{cor:zero-right-additive-identity-on-q}
For every $x\in\mathbb{Q}$,
\[
x+0_{\mathbb{Q}}=x.
\]
\end{corollary}
\end{corollarybox}

\begin{theorembox}{Theorem}
\begin{theorem}[Zero Is an Additive Identity on $\mathbb{Q}$]
\label{thm:zero-additive-identity-on-q}
For every $x\in\mathbb{Q}$,
\[
0_{\mathbb{Q}}+x=x
\qquad\text{and}\qquad
x+0_{\mathbb{Q}}=x.
\]
\end{theorem}
\end{theorembox}

\subsection*{Negation on $\mathbb{Q}$-Classes}

\begin{definitionbox}{Definition}
\begin{definition}[Negation on $\mathbb{Q}$-Classes]
\label{def:negation-on-q-classes}
Define negation on representatives by
\[
-[a,b]:=[-a,b].
\]
\end{definition}
\end{definitionbox}

\begin{dependencies}
\begin{itemize}
  \item \hyperref[def:rational-class-notation]{Rational Class Notation}
  \item \hyperref[def:negation-on-z-classes]{Negation on $\mathbb{Z}$}
\end{itemize}
\end{dependencies}

\begin{lemmabox}{Lemma}
\begin{lemma}[Negation Respects Fraction Equivalence]
\label{lem:negation-on-q-respects-equivalence}
If $[a,b]=[a',b']$, then
\[
-[a,b]=-[a',b'].
\]
\end{lemma}
\end{lemmabox}

\begin{theorembox}{Theorem}
\begin{theorem}[Negation on $\mathbb{Q}$ Is Well-Defined]
\label{thm:negation-on-q-well-defined}
Negation determines a well-defined unary operation $\mathbb{Q}\to\mathbb{Q}$.
\end{theorem}
\end{theorembox}

\begin{theorembox}{Theorem}
\begin{theorem}[Negation Gives a Left Additive Inverse on $\mathbb{Q}$]
\label{thm:negation-left-additive-inverse-on-q}
For every $x\in\mathbb{Q}$,
\[
(-x)+x=0_{\mathbb{Q}}.
\]
\end{theorem}
\end{theorembox}

\begin{corollarybox}{Corollary}
\begin{corollary}[Negation Gives a Right Additive Inverse on $\mathbb{Q}$]
\label{cor:negation-right-additive-inverse-on-q}
For every $x\in\mathbb{Q}$,
\[
x+(-x)=0_{\mathbb{Q}}.
\]
\end{corollary}
\end{corollarybox}

\begin{theorembox}{Theorem}
\begin{theorem}[Every Rational Has an Additive Inverse]
\label{thm:every-rational-has-additive-inverse}
For every $x\in\mathbb{Q}$ there exists $y\in\mathbb{Q}$ such that
\[
x+y=0_{\mathbb{Q}}
\qquad\text{and}\qquad
y+x=0_{\mathbb{Q}}.
\]
\end{theorem}
\end{theorembox}

\subsection*{Subtraction on $\mathbb{Q}$}

\begin{definitionbox}{Definition}
\begin{definition}[Subtraction on $\mathbb{Q}$]
\label{def:subtraction-on-q}
For $x,y\in\mathbb{Q}$, define
\[
x-y:=x+(-y).
\]
\end{definition}
\end{definitionbox}

\begin{dependencies}
\begin{itemize}
  \item \hyperref[def:addition-on-q-classes]{Addition on $\mathbb{Q}$}
  \item \hyperref[def:negation-on-q-classes]{Negation on $\mathbb{Q}$}
\end{itemize}
\end{dependencies}

\begin{theorembox}{Theorem}
\begin{theorem}[Subtraction on $\mathbb{Q}$ Is Well-Defined]
\label{thm:subtraction-on-q-well-defined}
Subtraction determines a well-defined binary operation
$\mathbb{Q}\times\mathbb{Q}\to\mathbb{Q}$.
\end{theorem}
\end{theorembox}

\subsection*{Multiplication on $\mathbb{Q}$-Classes}

\begin{definitionbox}{Definition}
\begin{definition}[Multiplication on $\mathbb{Q}$-Classes]
\label{def:multiplication-on-q-classes}
Define multiplication on representatives by
\[
[a,b]\cdot[c,d]:=[a\cdot c,\ b\cdot d].
\]
\end{definition}
\end{definitionbox}

\begin{dependencies}
\begin{itemize}
  \item \hyperref[def:rational-class-notation]{Rational Class Notation}
  \item \hyperref[def:multiplication-on-z-classes]{Multiplication on $\mathbb{Z}$}
\end{itemize}
\end{dependencies}

\begin{lemmabox}{Lemma}
\begin{lemma}[Multiplication Lands in $\operatorname{Pre}\mathbb{Q}$]
\label{lem:multiplication-denominator-nonzero-q}
If $b\ne 0_{\mathbb{Z}}$ and $d\ne 0_{\mathbb{Z}}$, then
\[
b\cdot d\ne 0_{\mathbb{Z}}.
\]
\end{lemma}
\end{lemmabox}

\begin{dependencies}
\begin{itemize}
  \item \hyperref[thm:z-has-no-zero-divisors]{$\mathbb{Z}$ Has No Zero Divisors}
\end{itemize}
\end{dependencies}

\begin{lemmabox}{Lemma}
\begin{lemma}[Multiplication Respects Fraction Equivalence on the Left]
\label{lem:multiplication-on-q-respects-left-equivalence}
If $[a,b]=[a',b']$, then
\[
[a,b]\cdot[c,d]=[a',b']\cdot[c,d].
\]
\end{lemma}
\end{lemmabox}

\begin{dependencies}
\begin{itemize}
  \item \hyperref[def:multiplication-on-q-classes]{Multiplication on $\mathbb{Q}$}
  \item \hyperref[def:fraction-equivalence]{Fraction Equivalence}
  \item \hyperref[prop:equality-criterion-for-rational-classes]{Equality Criterion for Rational Classes}
  \item \hyperref[cor:multiplicative-cancellation-on-z]{Multiplicative Cancellation on $\mathbb{Z}$}
\end{itemize}
\end{dependencies}

\begin{lemmabox}{Lemma}
\begin{lemma}[Multiplication Respects Fraction Equivalence on the Right]
\label{lem:multiplication-on-q-respects-right-equivalence}
If $[c,d]=[c',d']$, then
\[
[a,b]\cdot[c,d]=[a,b]\cdot[c',d'].
\]
\end{lemma}
\end{lemmabox}

\begin{dependencies}
\begin{itemize}
  \item \hyperref[def:multiplication-on-q-classes]{Multiplication on $\mathbb{Q}$}
  \item \hyperref[def:fraction-equivalence]{Fraction Equivalence}
  \item \hyperref[prop:equality-criterion-for-rational-classes]{Equality Criterion for Rational Classes}
  \item \hyperref[cor:multiplicative-cancellation-on-z]{Multiplicative Cancellation on $\mathbb{Z}$}
\end{itemize}
\end{dependencies}

\begin{lemmabox}{Lemma}
\begin{lemma}[Multiplication Respects Fraction Equivalence]
\label{lem:multiplication-on-q-respects-equivalence}
If $[a,b]=[a',b']$ and $[c,d]=[c',d']$, then
\[
[a,b]\cdot[c,d]=[a',b']\cdot[c',d'].
\]
\end{lemma}
\end{lemmabox}

\begin{dependencies}
\begin{itemize}
  \item \hyperref[def:multiplication-on-q-classes]{Multiplication on $\mathbb{Q}$}
  \item \hyperref[lem:multiplication-on-q-respects-left-equivalence]{Multiplication Respects Fraction Equivalence on the Left}
  \item \hyperref[lem:multiplication-on-q-respects-right-equivalence]{Multiplication Respects Fraction Equivalence on the Right}
\end{itemize}
\end{dependencies}

\begin{theorembox}{Theorem}
\begin{theorem}[Multiplication on $\mathbb{Q}$ Is Well-Defined]
\label{thm:multiplication-on-q-well-defined}
Multiplication determines a well-defined binary operation
$\mathbb{Q}\times\mathbb{Q}\to\mathbb{Q}$.
\end{theorem}
\end{theorembox}

\begin{dependencies}
\begin{itemize}
  \item \hyperref[def:multiplication-on-q-classes]{Multiplication on $\mathbb{Q}$}
  \item \hyperref[lem:multiplication-on-q-respects-equivalence]{Multiplication Respects Fraction Equivalence}
\end{itemize}
\end{dependencies}

\begin{theorembox}{Theorem}
\begin{theorem}[Multiplication on $\mathbb{Q}$ Is Associative]
\label{thm:multiplication-on-q-associative}
For all $x,y,z\in\mathbb{Q}$,
\[
(x\cdot y)\cdot z=x\cdot(y\cdot z).
\]
\end{theorem}
\end{theorembox}

\begin{theorembox}{Theorem}
\begin{theorem}[Multiplication on $\mathbb{Q}$ Is Commutative]
\label{thm:multiplication-on-q-commutative}
For all $x,y\in\mathbb{Q}$,
\[
x\cdot y=y\cdot x.
\]
\end{theorem}
\end{theorembox}

\begin{dependencies}
\begin{itemize}
  \item \hyperref[def:multiplication-on-q-classes]{Multiplication on $\mathbb{Q}$}
  \item \hyperref[thm:multiplication-on-q-well-defined]{Multiplication on $\mathbb{Q}$ Is Well-Defined}
  \item \hyperref[thm:multiplication-on-z-commutative]{Multiplication on $\mathbb{Z}$ Is Commutative}
\end{itemize}
\end{dependencies}

\begin{theorembox}{Theorem}
\begin{theorem}[One Is a Left Multiplicative Identity on $\mathbb{Q}$]
\label{thm:one-left-multiplicative-identity-on-q}
For every $x\in\mathbb{Q}$,
\[
1_{\mathbb{Q}}\cdot x=x.
\]
\end{theorem}
\end{theorembox}

\begin{corollarybox}{Corollary}
\begin{corollary}[One Is a Right Multiplicative Identity on $\mathbb{Q}$]
\label{cor:one-right-multiplicative-identity-on-q}
For every $x\in\mathbb{Q}$,
\[
x\cdot 1_{\mathbb{Q}}=x.
\]
\end{corollary}
\end{corollarybox}

\begin{theorembox}{Theorem}
\begin{theorem}[One Is a Multiplicative Identity on $\mathbb{Q}$]
\label{thm:one-multiplicative-identity-on-q}
For every $x\in\mathbb{Q}$,
\[
1_{\mathbb{Q}}\cdot x=x
\qquad\text{and}\qquad
x\cdot 1_{\mathbb{Q}}=x.
\]
\end{theorem}
\end{theorembox}

\subsection*{Reciprocal on Nonzero $\mathbb{Q}$-Classes}

\begin{lemmabox}{Lemma}
\begin{lemma}[Nonzero Predicate Is Well-Defined on $\mathbb{Q}$]
\label{lem:nonzero-predicate-well-defined-q}
For every prefraction $(a,b)\in\operatorname{Pre}\mathbb{Q}$,
\[
[a,b]\ne 0_{\mathbb{Q}}
\quad\Longleftrightarrow\quad
a\ne 0_{\mathbb{Z}}.
\]
Thus nonzeroness is independent of the representative.
\end{lemma}
\end{lemmabox}

\begin{definitionbox}{Definition}
\begin{definition}[Reciprocal on Nonzero $\mathbb{Q}$-Classes]
\label{def:reciprocal-on-q-classes}
For $[a,b]\ne 0_{\mathbb{Q}}$, define
\[
[a,b]^{-1}:=[b,a].
\]
\end{definition}
\end{definitionbox}

\begin{lemmabox}{Lemma}
\begin{lemma}[Reciprocal Lands in $\operatorname{Pre}\mathbb{Q}$]
\label{lem:reciprocal-denominator-nonzero-q}
If $[a,b]\ne 0_{\mathbb{Q}}$, then $a\ne 0_{\mathbb{Z}}$, so $(b,a)$ is a
prefraction.
\end{lemma}
\end{lemmabox}

\begin{lemmabox}{Lemma}
\begin{lemma}[Reciprocal Respects Fraction Equivalence]
\label{lem:reciprocal-on-q-respects-equivalence}
If $[a,b]=[a',b']$ and $[a,b]\ne 0_{\mathbb{Q}}$, then
\[
[a,b]^{-1}=[a',b']^{-1}.
\]
\end{lemma}
\end{lemmabox}

\begin{theorembox}{Theorem}
\begin{theorem}[Reciprocal on $\mathbb{Q}$ Is Well-Defined]
\label{thm:reciprocal-on-q-well-defined}
Reciprocal determines a well-defined map
\[
\mathbb{Q}\setminus\{0_{\mathbb{Q}}\}\to\mathbb{Q}\setminus\{0_{\mathbb{Q}}\}.
\]
\end{theorem}
\end{theorembox}

\begin{theorembox}{Theorem}
\begin{theorem}[Reciprocal Gives a Left Multiplicative Inverse on $\mathbb{Q}$]
\label{thm:reciprocal-left-multiplicative-inverse-on-q}
For every nonzero $x\in\mathbb{Q}$,
\[
x^{-1}\cdot x=1_{\mathbb{Q}}.
\]
\end{theorem}
\end{theorembox}

\begin{corollarybox}{Corollary}
\begin{corollary}[Reciprocal Gives a Right Multiplicative Inverse on $\mathbb{Q}$]
\label{cor:reciprocal-right-multiplicative-inverse-on-q}
For every nonzero $x\in\mathbb{Q}$,
\[
x\cdot x^{-1}=1_{\mathbb{Q}}.
\]
\end{corollary}
\end{corollarybox}

\begin{theorembox}{Theorem}
\begin{theorem}[Every Nonzero Rational Has a Multiplicative Inverse]
\label{thm:every-nonzero-rational-has-multiplicative-inverse}
For every $x\in\mathbb{Q}$ with $x\ne 0_{\mathbb{Q}}$, there exists
$y\in\mathbb{Q}$ such that
\[
x\cdot y=1_{\mathbb{Q}}
\qquad\text{and}\qquad
y\cdot x=1_{\mathbb{Q}}.
\]
\end{theorem}
\end{theorembox}

\subsection*{Division on $\mathbb{Q}$}

\begin{definitionbox}{Definition}
\begin{definition}[Division on $\mathbb{Q}$]
\label{def:division-on-q}
For $x,y\in\mathbb{Q}$ with $y\ne 0_{\mathbb{Q}}$, define
\[
x/y:=x\cdot y^{-1}.
\]
\end{definition}
\end{definitionbox}

\begin{theorembox}{Theorem}
\begin{theorem}[Division on $\mathbb{Q}$ Is Well-Defined]
\label{thm:division-on-q-well-defined}
Division determines a well-defined map
\[
\mathbb{Q}\times(\mathbb{Q}\setminus\{0_{\mathbb{Q}}\})\to\mathbb{Q}.
\]
\end{theorem}
\end{theorembox}

\subsection*{Distributivity on $\mathbb{Q}$-Classes}

\begin{theorembox}{Theorem}
\begin{theorem}[Left Distributivity on $\mathbb{Q}$]
\label{thm:left-distributivity-on-q}
For all $x,y,z\in\mathbb{Q}$,
\[
x\cdot(y+z)=x\cdot y+x\cdot z.
\]
\end{theorem}
\end{theorembox}

\begin{corollarybox}{Corollary}
\begin{corollary}[Right Distributivity on $\mathbb{Q}$]
\label{cor:right-distributivity-on-q}
For all $x,y,z\in\mathbb{Q}$,
\[
(y+z)\cdot x=y\cdot x+z\cdot x.
\]
\end{corollary}
\end{corollarybox}

\begin{theorembox}{Theorem}
\begin{theorem}[Multiplication Distributes over Addition on $\mathbb{Q}$]
\label{thm:multiplication-distributes-over-addition-on-q}
Multiplication distributes over addition on both sides in $\mathbb{Q}$.
\end{theorem}
\end{theorembox}

\subsection*{Class-Level Field Summary}

\begin{theorembox}{Theorem}
\begin{theorem}[$\mathbb{Q}$ Is an Additive Abelian Group]
\label{thm:q-additive-abelian-group}
The structure $(\mathbb{Q},+,0_{\mathbb{Q}})$ is an abelian group.
\end{theorem}
\end{theorembox}

\begin{theorembox}{Theorem}
\begin{theorem}[Nonzero $\mathbb{Q}$ Is a Multiplicative Abelian Group]
\label{thm:q-nonzero-multiplicative-abelian-group}
The structure
$(\mathbb{Q}\setminus\{0_{\mathbb{Q}}\},\cdot,1_{\mathbb{Q}})$ is an abelian
group.
\end{theorem}
\end{theorembox}

\begin{theorembox}{Theorem}
\begin{theorem}[$\mathbb{Q}$ Is a Field]
\label{thm:q-is-a-field}
The structure $(\mathbb{Q},+,\cdot,0_{\mathbb{Q}},1_{\mathbb{Q}})$ is a field.
\end{theorem}
\end{theorembox}

\begin{dependencies}
\begin{itemize}
  \item \hyperref[thm:zero-not-one-in-q]{Zero Is Not One in $\mathbb{Q}$}
\end{itemize}
\end{dependencies}
